% SGA 3, 第 VI_B 篇, 第 1 轮英文重构.
% 范围: 推论 5.7.1, 注 5.7.2, 段落 5.8 标题, 命题 5.8.1 与推论 5.8.2 全文,
% 含底本编者注 47.
% 底本: Polo--Gille Exp6B-13oct24.pdf, 分篇第 35 页
% (合订本第 376 页; 原版页边页码 357, 重版本页眉 365);
% 注 47 核对至分篇第 36 页 / 合订本第 377 页顶部的续文.
% 例 5.8.3 始于分篇第 35 页 / 合订本第 376 页.
% OCR 与 Reinhold 英文 Markdown 仅为比对材料.
% 本单元无图.

\medskip
\phantomsection\label{cor:sga3-VIB-5.7.1}%
\noindent\textbf{推论 5.7.1.---}
\emph{设 \(S\) 和 \(G\) 如 \hyperref[thm:sga3-VIB-5.7]{5.7} 中所述,
\(\sigma,\tau\) 是 \(G^0\) 的两个 \(S\)-截面,
即 \(\sigma,\tau\in G^0(S)\).
则 \(\sigma\) 和 \(\tau\) 的重合子概形
\(S(\sigma,\tau)\subset S\)
\textup{(即 \(G\times_S G\) 的对角线在态射 \((\sigma,\tau)\) 下的逆像)} 闭.}

事实上, 对每个 \(s\in S\), 设 \(U\) 是 \(G\)
的一个包含 \(\epsilon(s)\) 的仿射开集. 则
\[
  V=\epsilon^{-1}(U)
\]
是 \(s\) 在 \(S\) 中的开邻域,
并且由于 \(G^0U\) 分离, \(S(\sigma,\tau)\cap V\) 在 \(V\) 中闭.

\medskip
\phantomsection\label{rem:sga3-VIB-5.7.2}%
\noindent\textbf{注 5.7.2.---}
Gabber 指出可以证明如下事实: 若 \(S\) 是以 \(s\) 为闭点的局部 Hensel 概形,
则开集 \(G^0U\) 的交是 \(G\) 的一个在 \(S\) 上分离的开子群概形,
其中 \(U\) 遍历 \(\epsilon(s)\) 的仿射开邻域.

\section*{5.8. 补充}
\phantomsection\label{sec:sga3-VIB-5.8}%
\addcontentsline{toc}{section}{\SGARefLink{sec:sga3-VIB-5.8}{5.8}. 补充}

\noindent\footnote{\label{note:sga3-VIB-47}编者注: 此段补充与反例为我们新增,
全部由 O.\ Gabber 告知.}%
先回忆\hyperref[prop:sga3-VIA-2.6.6]{命题 VI\(_A\), 2.6.6}:

\medskip
\phantomsection\label{prop:sga3-VIB-5.8.1}%
\noindent\textbf{命题 5.8.1.---}
\emph{设 \(k\) 是域, \(G\) 是局部有限型 \(k\)-群,
它作用于 \(k\)-概形 \(X\), 使态射}
\[
  G\times X\longrightarrow X\times X,
  \qquad (g,x)\longmapsto(gx,x),
\]
\emph{满射. 假设 \(X\) 拟分离. 则 \(X\) 的连通分支不可约.}

\medskip
\phantomsection\label{cor:sga3-VIB-5.8.2}%
\noindent\textbf{推论 5.8.2.---}
\emph{设 \(S,G,X\) 如 \hyperref[thm:sga3-VIB-5.3]{5.3} 的预备假设中所述.
再假设每个纤维 \(X_s\) 都拟分离且连通.
则 \(X\) 在 \(S\) 上分离且拟紧.}

事实上, 由前一个命题, 每个纤维 \(X_s\) 都不可约,
证明的其余部分与 \hyperref[cor:sga3-VIB-5.4]{5.4} 的证明相同.
